\section{Introduction}
An engineer testing a nested virtual machine may control its containing Windows
environment while using a binary-only intermediate virtual machine monitor
(VMM). A useful experiment is to substitute a page seen by the nested guest,
observe the application effect, and restore the original page without editing
the VMM. This requires more than obtaining a physical address: the page belongs
to several translation domains, cached translations can outlive a policy
change, and a reused address need not identify the same workload object.

This paper presents \ks{}, an experimental Intel VMX monitor loaded from an
authorized Windows kernel inside an existing Hyper-V guest. It interposes below
that running Windows instance and subsequently hosts VMware and a TinyCore Linux
descendant. Its page-control interface redirects one selected 4\,KiB descendant
guest-physical page --- or a 2\,MiB region of them, cloned from the original so
that publication changes nothing the descendant can observe --- to private
backing, and exposes publication, invalidation, rollback, and reclamation as
separate events. Guest-side helpers select, pin, and observe the experimental
page; the page intervention itself is controlled from the containing Windows
instance.

The distinction between insertion and later page control is essential.
\emph{The evaluated insertion starts while descendant VMMs are absent.}
Windows continues within the same boot and with the same observed processes.
VMware and TinyCore start afterward. Their identities are then preserved during
page substitution and recovery. Inserting beneath a VMware/TinyCore stack that
is already running would additionally require transferring existing VMX
ownership and execution state; every insertion measured here begins in the
first arrangement, and we report no measurement of the second.
Figure~\ref{fig:topology} makes the two intervals explicit.

\begin{figure}[tbp]
\centering
\begin{tikzpicture}[font=\small, node distance=4mm,
  stage/.style={draw=blue!45!black,fill=blue!3,rounded corners=2pt,
                align=left,text width=6.65cm,inner sep=8pt},
  arrow/.style={-{Latex[length=2.2mm]},thick,draw=blue!55!black}]
\node[stage] (a) {\textbf{1. Live insertion}\par\medskip
  Outer Hyper-V\par
  \hspace*{3mm}Root: Windows 0 / HVCI\par
  \hspace*{3mm}Guest: running Windows 1\par\medskip
  \textbf{Intermediate VMM absent}\par
  Insert KSword; preserve Windows boot\par
  and observed process identities.};
\node[stage,right=10mm of a] (b) {\textbf{2. Later descendant execution}\par\medskip
  Outer Hyper-V\par
  \hspace*{3mm}Root: Windows 0 / HVCI\par
  \hspace*{3mm}Guest: \textcolor{blue!65!black}{\textbf{KSword VMX}}\par
  \hspace*{8mm}Same Windows 1 (4 vCPUs)\par
  \hspace*{13mm}VMware (started afterward)\par
  \hspace*{18mm}TinyCore (2 vCPUs)\par
  \hspace*{18mm}\textcolor{blue!65!black}{Selected 4 KiB page}};
\draw[arrow] (a.east) -- (b.west);
\node[draw=green!40!black,fill=green!3,rounded corners=2pt,
      align=center,below=7mm of b,text width=6.65cm,inner sep=8pt] (c)
  {\textbf{Reversible composed-EPT control}\par\medskip
   Original backing $\longrightarrow$ private backing\par
   $\longrightarrow$ original backing after remove + drain};
\draw[arrow] (b.south) -- (c.north);
\node[align=left,text width=6.65cm,below=9mm of a,inner sep=4pt]
  {\textbf{Continuity is measured per interval.}\par
   Insertion and page control use separate\par
   continuity intervals.};
\end{tikzpicture}

\caption{Evaluated topology and temporal order. First, \ks{} is inserted beneath
running Windows~1 while intermediate VMMs are absent. VMware and TinyCore are
then started normally. The subsequent page-control interval preserves the
observed Windows, VMware, and TinyCore identities. Windows~0 and outer Hyper-V
remain outside the inserted monitor. This figure does not depict insertion
beneath an already-running descendant stack.}
\label{fig:topology}
\end{figure}

Neither late OS virtualization nor translation composition originates with this
work. SubVirt placed a VMM beneath existing operating systems, and Blue Pill
described on-the-fly interposition~\cite{subvirt,bluepill}. Turtles established
nested VMX and multi-dimensional paging for unmodified
VMMs~\cite{turtles}. We build on these concepts and investigate a particular
deployment and lifecycle problem: a monitor started inside a Hyper-V Windows
guest that controls a later descendant without an intermediate-VMM extension.
The intended use is controlled systems experimentation, not covert persistence
or a new protection boundary against the surrounding virtualization stack.

The paper makes three bounded contributions:
\begin{enumerate}
\item A Windows implementation combining live Windows-only interposition,
      nested VMX execution, and a descendant page lease. The lease binds a
      referenced VMM process identity to a sampled source translation path and
      retains replacement backing until an all-CPU retirement drain succeeds.
      A region larger than a page is admitted only on a property that makes one
      leaf safe for all of it --- either the source's own leaf is already that
      coarse, or every source entry under the region is read and found to agree
      --- and the second, weaker basis is rechecked by sampling and revoked
      under its own reason when it stops holding.
\item An evaluation that connects control events to full-page observations on
      two descendant vCPUs, identity continuity, fault outcomes, and allocation
      ledgers. A BusyBox HTTP service and a fixed-schema policy consumer show
      an application-visible data fault and recovery, with a guest direct-write
      comparator that makes the differing control requirements explicit.
\item A version-separated evidence set reporting useful behavior together with
      substantial costs and negative results: millisecond-scale insertion
      bodies, paced application observations, residual CPUID/RTT overhead,
      refused admission in an inner Hyper-V root partition, a 1\,GiB leaf
      requested on hardware and refused under both rules, and the finding that
      which guest-physical bases can be replaced at all is decided by what the
      descendant has touched rather than by how much memory it was given.
\end{enumerate}

All measurements come from one physical machine. In the current cohort, five
measured insertions have a median driver-body time of 2.5953\,ms. Twenty-four
paced TCP runs, including warmups, produce 74,145 valid responses. Three current
two-vCPU page cycles and fifteen injected-control trials complete, while one of
four HTTP EPT attempts is incomplete because its bounded observer expires.
The current nested mode incurs 8.76 times the off-mode CPUID latency and
30.7\% higher Windows loopback RTT. These results establish a working, limited
control mechanism; they do not establish low overhead across applications,
arbitrary page-update atomicity, or general hot insertion beneath existing VMMs.

\section{Background and operating assumptions}
\label{sec:background}
\paragraph{Topology names.}
We call the outer Hyper-V root operating system \emph{Windows~0} and its guest
\emph{Windows~1}. The \emph{intermediate VMM} is VMware running within Windows~1;
the \emph{descendant} is its TinyCore guest. After insertion, \ks{} monitors both
Windows~1 and the nested VMX execution it presents to VMware. These names avoid
assigning a single global meaning to L0/L1/L2 across two nested-virtualization
implementations. Outer Hyper-V continues to control actual machine execution.

\paragraph{VMX and address translation.}
VMX separates host and guest execution state using virtual-machine control
structures (VMCSs). Extended page tables (EPT) translate guest-physical addresses
and carry access permissions and memory types. Cached EPT-derived translations
require the appropriate invalidation when mappings change; changing a software
pointer alone does not retire those translations~\cite{intel}. In \ks{}'s local
naming, EPT12 is VMware's source translation, and the composed or shadow EPT
is the translation that \ks{} supplies for descendant execution. A physical
address reported by this monitor is a Windows~1 physical address, not a claim
to know the outer machine's system-physical address.

\paragraph{Authority and trust.}
The operator authorizes the driver and has the Windows~1 kernel privileges
needed to install it. Outer Hyper-V, Windows~0, and the active outer HVCI
configuration are trusted. The intermediate VMM is assumed to obey architectural
invalidation rules. The prototype does not defend against a malicious VMM that
races translation changes, an untrusted outer hypervisor, or an adversarial
guest deliberately reusing the selected object. Firmware settings and nested
capability exposure are prerequisites, not privileges acquired by the monitor.

\paragraph{Control-path independence.}
Page intervention uses the \ks{} driver interface without a VMware/Hyper-V
management API, VMM code hook, or component injected into the VMM. This is a
statement about the intervention path. Experiment setup uses VM management,
PowerShell Direct, an ISO image, and console/serial access. Normal hypercalls,
nested VMX instructions, and EPT accessed/dirty (A/D) maintenance still occur;
some A/D maintenance writes translation metadata owned by the intermediate VMM.
Thus, the claim is neither that the VMMs are never invoked nor that all their
memory remains byte-for-byte unchanged.
